Para responder al problema descrito y a las carencias identificadas en las
plataformas existentes, este \acs{tfg} persigue los siguientes objetivos:

\begin{objetive}
  \item Diseñar e implementar el \dfn{backend} completo del SCDEE como un
        sistema distribuido \dfn{multi-tenant}.
        \label{OBJ:BACKEND-MULTITENANT}
  \item Dotar al \dfn{backend} de una arquitectura modular y extensible,
        sustentada en una clara separación de responsabilidades y en
        interfaces de \dfnpl{plugin} intercambiables.
        \label{OBJ:MODULAR-ARCH}
  \item Construir un \textit{pipeline} de ingesta y reconocimiento de
        páginas escaneadas que combine identificación por código \acs{qr},
        \acs{ocr} y un \textit{matching} robusto a errores de
        reconocimiento.
        \label{OBJ:OCR-PIPELINE}
  \item Cubrir el ciclo operativo completo de corrección: definición
        flexible de zonas \acs{ocr} por perfil de página, asignación de
        correctores mediante reglas configurables, anotaciones de los
        correctores sobre las páginas escaneadas y entrega controlada de
        los exámenes corregidos a los alumnos durante la ventana de
        revisión.
        \label{OBJ:GRADING-CYCLE}
  \item Ofrecer un sistema de configuración dinámica por organización y
        tareas de mantenimiento automatizadas, permitiendo a cada
        \dfn{tenant} adaptar el sistema a sus necesidades sin intervención
        del administrador.
        \label{OBJ:CONFIG-MAINTENANCE}
  \item Garantizar el cumplimiento de los requisitos de seguridad y
        protección de datos: cifrado en reposo de datos sensibles,
        auditoría inmutable, \textit{soft delete} y mecanismos de
        protección frente a la fuga de exámenes.
        \label{OBJ:SECURITY-RGPD}
  \item Generar una documentación OpenAPI completa y autoexplicativa que
        actúe como contrato formal entre \dfn{backend} y \dfn{frontend}.
        \label{OBJ:OPENAPI-CONTRACT}
  \item Validar el sistema mediante pruebas unitarias, de integración y de
        carga, comprobando el cumplimiento de los requisitos de rendimiento
        y concurrencia objetivo.
        \label{OBJ:VALIDATION}
\end{objetive}

El grado de consecución de cada uno de estos objetivos se discute en
el \cref{CAP:CONCLUSIONS}.
